\newglossaryentry{dsrc-gloss}{%
  name = {DSRC},
  description = {Dedicated Short-Range Communications; standardisierte Kurzstreckenkommunikation für Fahrzeug-zu-Infrastruktur-Anwendungen}
}

\newglossaryentry{gnss-gloss}{%
  name = {GNSS},
  description = {Global Navigation Satellite System; Oberbegriff für satellitengestützte Navigationssysteme wie GPS, Galileo, GLONASS oder BeiDou}
}

\newglossaryentry{ism-band}{%
  name = {ISM-Band},
  description = {Industrial, Scientific and Medical Band; lizenzfreier Frequenzbereich für Industrie-, Wissenschafts- und Medizinalanwendungen, u.\,a. bei 2{,}4\,GHz}
}

\newglossaryentry{gatt-gloss}{%
  name = {GATT},
  description = {Generic Attribute Profile; BLE-Datenmodell zur Strukturierung von Services und Characteristics}
}

\newglossaryentry{tvs-gloss}{%
  name = {TVS},
  description = {Transient Voltage Suppressor; Schutzdiode zur Begrenzung kurzzeitiger Überspannungen}
}

\newglossaryentry{tx-gloss}{%
  name = {TX},
  description = {Transmit; Sendeleitung oder Sendeweg einer seriellen Schnittstelle}
}

\newglossaryentry{rx-gloss}{%
  name = {RX},
  description = {Receive; Empfangsleitung oder Empfangsweg einer seriellen Schnittstelle}
}

\newglossaryentry{parity-bit}{%
  name = {Paritätsbit},
  description = {Zusatzbit zur Fehlererkennung bei serieller Übertragung, meist als gerade oder ungerade Parität}
}

\newglossaryentry{soc-gloss}{%
  name = {SoC},
  description = {System-on-a-Chip; integrierter Schaltkreis, der Prozessor, Speicher und Peripherie in einem Baustein vereint}
}

\newglossaryentry{xtensa-cores-gloss}{%
  name = {Xtensa-Kerne},
  description = {Prozessorkerne der Tensilica-Xtensa-Architektur, wie sie u.\,a. im ESP32 eingesetzt werden}
}

\newglossaryentry{spi-gloss}{%
  name = {SPI},
  description = {Serial Peripheral Interface; synchrones serielles Bussystem für die Kommunikation zwischen Mikrocontroller und Peripherie}
}

\newglossaryentry{i2c-gloss}{%
  name = {I\textsuperscript{2}C},
  description = {Inter-Integrated Circuit; serielles Zwei-Draht-Bussystem für die Kommunikation zwischen Mikrocontrollern und Peripherie}
}

\newglossaryentry{adc-gloss}{%
  name = {ADC},
  description = {Analog-Digital-Wandler; wandelt analoge Signale in digitale Werte}
}

\newglossaryentry{freertos-gloss}{%
  name = {FreeRTOS},
  description = {Real-Time Operating System; quelloffenes Echtzeitbetriebssystem für Mikrocontroller}
}

\newglossaryentry{ttl-cmos-gloss}{%
  name = {TTL/CMOS-Logikpegel},
  description = {Logikpegel klassischer TTL- oder CMOS-Schaltungen, typischerweise wenige Volt}
}

\newglossaryentry{rs232-pegel-gloss}{%
  name = {RS-232-Pegel},
  description = {Spannungspegel der RS-232-Schnittstelle, typischerweise positive und negative Spannungen}
}

\newglossaryentry{anhang-ic-gloss}{%
  name = {Anhang IC},
  description = {Anhang der EU-Durchführungsverordnung zu digitalen Tachographen mit technischen Anforderungen und Datenstrukturen}
}

\newglossaryentry{ferritperle-gloss}{%
  name = {Ferritperle},
  description = {Ferritbauteil zur Dämpfung hochfrequenter Störungen in Leitungen}
}

\newglossaryentry{elko-gloss}{%
  name = {Elko},
  description = {Elektrolytkondensator zur Energiespeicherung und Glättung von Versorgungsspannungen}
}

\newglossaryentry{low-esr-gloss}{%
  name = {Low ESR},
  description = {Geringer äquivalenter Serienwiderstand zur Reduktion von Verlusten und Ripple}
}

\newglossaryentry{mlcc-gloss}{%
  name = {MLCC},
  description = {Multilayer Ceramic Capacitor; mehrlagiger Keramikkondensator}
}

% Keramik-Dielektrikum als Glossarbegriff statt X7R-Abkürzung
\newglossaryentry{keramik-gloss}{%
  name = {Keramik},
  description = {Keramik-Dielektrikum für Kondensatoren mit stabilen Kapazitätswerten über einen breiten Temperaturbereich}
}

\newglossaryentry{isat-gloss}{%
  name = {I\textsubscript{sat}},
  description = {Sättigungsstrom einer Induktivität, ab dem die Induktivität deutlich abnimmt}
}

\newglossaryentry{schottky-gloss}{%
  name = {Schottky-Diode},
  description = {Diode mit geringerer Durchlassspannung und schnellen Schaltzeiten}
}

\newglossaryentry{fr4-gloss}{%
  name = {FR4},
  description = {Glasfaserverstärktes Epoxidharz als Standard-Leiterplattenmaterial}
}

\newglossaryentry{hal-gloss}{%
  name = {HAL},
  description = {Hot Air Leveling; Oberflächenfinish zur Verzinnung von Leiterplatten}
}

\newglossaryentry{enig-gloss}{%
  name = {ENIG},
  description = {Electroless Nickel Immersion Gold; Oberflächenfinish mit Nickel- und Goldschicht}
}

\newglossaryentry{esd-gloss}{%
  name = {ESD},
  description = {Electrostatic Discharge; elektrostatische Entladung, die elektronische Bauteile schädigen kann}
}
